\section{Related Work}
\label{sec:related}

\subsection{LLM GUI Agent Security}
\label{sec:rw_agents}

LLM-driven mobile GUI agents have advanced rapidly in capability and deployment.%
AppAgent~\cite{appagent2024} and AutoGLM~\cite{autoglm2024} both achieve
task-completion rates above 36\% on general phone benchmarks, with AutoGLM commercially
deployed.%
CogAgent~\cite{cogagent2024} achieves state-of-the-art GUI navigation using a
dual-resolution 18B VLM without system API access; Mobile-Agent~\cite{mobileagent2024}
enables vision-only phone operation without XML metadata.%
SeeAct~\cite{seeact2024} demonstrates that GPT-4V can act as a generalist web agent when
paired with effective grounding, and AndroidWorld~\cite{androidworld2024} standardizes
evaluation across 116 tasks on 20 real Android apps.%
Anthropic's Computer Use~\cite{anthropiccu2024} extends the paradigm to desktop interfaces.%
The security implications of agent-controlled mobile devices have been documented in the
fraud context: mass account registration, advertising click abuse, welfare farming, and
unauthorized payment confirmation~\cite{llmagentfraud2025,clickfraud2021,adffraud2021}.%
Existing defenses targeting LLM agents focus on prompt injection~\cite{agentinjection2024}
and policy-level access control~\cite{agentpolicy2024}, rather than detecting whether
an active session is agent-controlled.%
Our work addresses this orthogonal detection gap.%

\subsection{Mobile Bot and Click-Farm Detection}
\label{sec:rw_bot}

Classical mobile bot detection relies on device attestation (SafetyNet, Play Integrity API)
and emulator fingerprinting~\cite{emulatordetect2022}.%
Phone farms defeat both: real hardware passes attestation and carries genuine device
identifiers.%
Environment-aware evasion techniques~\cite{sandboxevasion2022} further demonstrate that
distinguishing real devices from sandboxes via user-artifact differences defeats
state-of-the-art mobile sandboxes.%
Statistical behavioral methods---touch trajectory analysis, UI interaction graphs---were
designed for scripted bots~\cite{frauddroid2017} and are bypassed by LLM agents that
adapt to UI state.%
The Android accessibility framework has been extensively misused by malware for
cross-app data theft and unauthorized transaction fraud~\cite{a11yabuse2021,dva2024},
illustrating the breadth of the software injection threat surface.%
Large-scale measurement studies confirm that click-farm-based invalid traffic already
operates at industrial scale~\cite{clickfraud2021,adffraud2021}, motivating detection
primitives that remain valid against real-hardware phone farms.%
BeCAPTCHA~\cite{becaptcha2021} combined accelerometer and touch kinematics for bot
detection in a CAPTCHA-style challenge, but its feature engineering targets
repetitive fixed-script bots and does not model the physical causality gap.%
Turing Test on Screen~\cite{ttos2026} is the closest prior work: it applies kinematic
features to LLM agent detection and reports strong results against vanilla agents, but
explicitly acknowledges that physical sensor spoofing lies outside its scope.%
Our work closes this gap with a structurally unforgeable physical primitive.%

\subsection{Sensor-Based Liveness and User Authentication}
\label{sec:rw_liveness}

Sensor-based liveness detection has been explored in the context of fingerprint and
face recognition~\cite{sensorliveness2023} and touch-screen authentication~\cite{touchauth2020}.%
USENIX Security has shown that even liveness is insufficient: puppet attacks can replay
captured finger motion, requiring behavioral biometrics as a complementary
layer~\cite{livenessnotenough2020}.%
TouchPass~\cite{touchpass2020} authenticates users via chassis vibrations during
touch---a physical signal related to, but distinct from, the causality impulse we exploit.%
Behavioral biometric systems such as MMAuth~\cite{mmauth2022} and
AuthentiSense~\cite{authentisense2023} achieve high accuracy for continuous
authentication, but remain vulnerable to synthesis attacks because they measure
statistical patterns rather than physical necessity.%
Modern CAPTCHA evaluations~\cite{captchastudy2023} confirm that purely behavioral
challenges are increasingly ineffective, with state-of-the-art solvers bypassing
them at high rates.%
BioMotoTouch~\cite{biomotouch2026} characterizes the mechanical impulse response of
real finger contact as measured by the accelerometer, providing the physical model
underlying Observation~1 (Section~\ref{sec:obs1}).%
iStelan~\cite{istelan2023} demonstrates that capacitive touchscreen contact produces
an EM transient detectable by the permission-free magnetometer.%
Zero-permission motion sensors have further been shown to leak speech
content~\cite{accear2022} and private voice-assistant responses~\cite{stealthyimu2023},
and PhyScout~\cite{physcout2024} detects multi-sensor spoofing via spatio-temporal
consistency---illustrating the richness of physical information available at zero
privilege.%
Our work adapts the accelerometer-based physical causality chain to the LLM agent
detection problem, introducing a new three-class threat model (including the novel
overlay agent class) and a cross-modal attention architecture optimized for this setting.%

\subsection{Cross-Modal and Self-Supervised Learning for Sensor Data}
\label{sec:rw_ml}

Self-supervised pre-training on sensor time-series has produced strong generalizable
representations.%
TF-C~\cite{tfc2022} enforces consistency between time-domain and frequency-domain views
via contrastive learning.%
PatchTST~\cite{patchtst2023} segments time series into patch-level tokens for
channel-independent Transformer processing with masked pre-training, directly
motivating our patch-based accelerometer encoder in Stage~1.%
SimMTM~\cite{simmtm2023} recovers masked time points via weighted neighbor aggregation
on the data manifold, providing an alternative masked pre-training objective for
sensor data.%
FOCAL~\cite{focal2023} uses factorized shared-and-private encoders for multi-sensor
cross-modal contrastive pre-training.%
CA-TCC~\cite{catcc2023} extends contrastive learning to semi-supervised sensor HAR.%
TS-TCC~\cite{tstcc2021} applies temporal and contextual contrastive objectives to IMU data.%
Crossformer~\cite{crossformer2023} introduces cross-dimension attention for multivariate
time series, capturing dependencies across sensor channels; our PC-CWA adapts the
cross-attention paradigm to the asymmetric, physics-constrained relationship between
the touch-event and accelerometer streams.%
CausalFormer~\cite{causalformer2024} discovers causal graphs from multivariate time
series via interpretable attention; our PC-CWA differs by encoding \textit{known}
physical causal structure as a hard architectural constraint rather than learning
causality from data.%
Our PCCN Stage~2 training loss builds on TF-C~\cite{tfc2022} and SupCon~\cite{supcon2020},
while our cross-modal attention head is informed by the cross-dimensional processing
principles of FOCAL and Crossformer, adapted to the asymmetric causality structure
of the accelerometer--touch relationship.%
